\section{Limitations and Ethics}
\label{sec:limitations}

\paragraph{Evaluation scope.}
Our functional-model experiments use two resettable shopping applications and
three runs per condition. This controlled setting permits shared candidates,
budgets, and evidence criteria, but does not establish performance across the
diversity of production Web applications. The induced model covers only states
and functions reached during exploration; it is neither a complete inventory
of application functionality nor a reconstruction of hidden business logic.
Location constraints and observed direct dependencies are local applicability
evidence, not complete preconditions or causal rules. Comparisons with systems
that use different tasks, action spaces, executors, or labels would require a
new common evaluation rather than direct comparison of reported scores.

\paragraph{Imperfect components and changing environments.}
VERA currently relies on vision-language models for hypothesis generation,
outcome judgment, and risk assessment, and on a browser executor for grounding
and interaction. Errors in any component can propagate into missed functions,
incorrect evidence judgments, or inaccurate risk records. Executor failure can
also be difficult to distinguish from an unavailable function. Moreover, Web
applications evolve: a claim supported at one time can become stale after an
interface or policy change. The current work retains provenance that can
support reinspection, but does not implement systematic staleness detection or
continuous revalidation.

\paragraph{Risk recognition is not intervention.}
The evaluated risk layer judges only the selected semantic action in its
current visual context; it does not scan every visible candidate or reason over
all possible long-horizon consequences. Its binary decision and single primary
type also compress compound or ambiguous risks. In C3, visual context improves
risk-type grounding most clearly, while binary detection trades higher recall
for lower precision and its recall and F1 difference intervals cross zero.
Risk records are produced before execution and could be consumed by a blocking
or human-confirmation policy, but the evaluated system does not use them to
alter actions. We therefore do not claim reduced harm, successful prevention,
or an end-to-end safety guarantee.

\paragraph{Annotation and statistical limitations.}
AI-assisted initial labels were reviewed and revised item by item by a human
author under frozen guides, but the study does not include independent double
annotation or inter-annotator agreement. The external context challenge
improves contextual diversity, yet its actions are supplied to the assessor and
do not show that VERA would propose them. Only eight candidate pairs realize a
binary gold-label difference, limiting pair-level conclusions. The three C2
repetitions are reported descriptively rather than used for significance
claims.

\paragraph{Ethical use of autonomous exploration.}
Open-ended agents can delete or disclose data, send communications, change
permissions, enter transactions, or consume service resources. Screenshots and
traces may also contain credentials or personal information. Experiments should
therefore use authorized, resettable environments and isolated accounts where
possible, minimize and redact sensitive data, retain only necessary artifacts,
and respect access controls, service terms, and applicable review requirements.
Deploying VERA on production systems would require an intervention policy,
failure handling, data governance, and human oversight beyond the mechanisms
evaluated here.
